\begin{apendicesenv}
\partapendices

% ----------------------------------------------------------------------
\chapter{Colinha de sintaxe}
\label{ap:colinha}

\begin{sqlcode}[caption={Estrutura (DDL)}]
CREATE DATABASE nome CHARACTER SET utf8mb4;
USE nome;
CREATE TABLE t (id INT AUTO_INCREMENT PRIMARY KEY, campo VARCHAR(100) NOT NULL);
ALTER TABLE t ADD COLUMN novo VARCHAR(50);
ALTER TABLE t MODIFY COLUMN campo VARCHAR(200) NOT NULL;
ALTER TABLE t DROP COLUMN novo;
DROP TABLE t;
\end{sqlcode}

\begin{sqlcode}[caption={Inspeção}]
SHOW DATABASES;   SHOW TABLES;   DESCRIBE t;   SHOW CREATE TABLE t;
\end{sqlcode}

\begin{sqlcode}[caption={Dados (DML e DQL)}]
INSERT INTO t (c1, c2) VALUES (v1, v2), (v3, v4);
SELECT c1, c2 FROM t WHERE cond ORDER BY c1 DESC LIMIT 10;
UPDATE t SET c1 = v WHERE id = 1;      -- SEMPRE com WHERE
DELETE FROM t WHERE id = 1;            -- SEMPRE com WHERE
\end{sqlcode}

\begin{sqlcode}[caption={Filtros}]
WHERE a = 1 AND b <> 2
WHERE a IN (1,2,3)          WHERE a BETWEEN 1 AND 10
WHERE nome LIKE 'A%'        WHERE campo IS NULL
\end{sqlcode}

\begin{sqlcode}[caption={Relacionamento}]
FOREIGN KEY (curso_id) REFERENCES curso(id) ON DELETE SET NULL;
SELECT m.nome, c.nome FROM membro m INNER JOIN curso c ON m.curso_id = c.id;
SELECT m.nome, c.nome FROM membro m LEFT  JOIN curso c ON m.curso_id = c.id;
\end{sqlcode}

\begin{sqlcode}[caption={Agregação}]
SELECT c.nome, COUNT(*) FROM curso c JOIN membro m ON m.curso_id = c.id
GROUP BY c.id, c.nome HAVING COUNT(*) >= 2;
\end{sqlcode}

% ----------------------------------------------------------------------
\chapter{Script completo da aula}
\label{ap:script}

Execute do início ao fim para reconstruir todo o cenário da capacitação. O
arquivo também está disponível no repositório, em
\comando{materials/liga.sql}.

\begin{aviso}[Antes de colar]
Substitua as três ocorrências de \comando{liga\_SEUNOME} pelo seu banco
--- \comando{liga\_ana}, \comando{liga\_joao} e assim por diante. Se você colar
sem trocar, vai criar um banco literalmente chamado \comando{liga\_seunome} e
disputá-lo com quem fizer o mesmo.
\end{aviso}

\lstinputlisting[style=sql]{../materials/liga.sql}

% ----------------------------------------------------------------------
\chapter{Exercícios}
\label{ap:exercicios}

Resolva no seu banco \comando{liga\_SEUNOME}, já construído pelo
Apêndice~\ref{ap:script}.

\begin{enumerate}
    \item Liste o nome e o e-mail de todos os membros \textbf{ativos}.
    \item Mostre os membros do 1º ao 5º período, do mais novo para o mais
          antigo na liga.
    \item Quantos membros ingressaram a partir de 2025?
    \item Cadastre um novo curso, \comando{'Engenharia Elétrica'}, e um membro
          nele.
    \item Corrija o e-mail da Carla para \comando{carla.dias@liga.edu.br}.
    \item Liste membro e nome do curso, incluindo quem está sem curso.
    \item Desative (\comando{ativo = FALSE}) todo membro que ingressou antes de
          \comando{2024-01-01}.
    \item Tente inserir um membro com \comando{curso\_id = 999}. Explique a
          mensagem de erro.
\end{enumerate}

\section{Respostas}

\begin{sqlcode}[caption={Exercícios 1 a 3}]
-- 1
SELECT nome, email FROM membro WHERE ativo = TRUE;

-- 2
SELECT * FROM membro
WHERE periodo BETWEEN 1 AND 5
ORDER BY data_ingresso DESC;

-- 3
SELECT COUNT(*) FROM membro WHERE data_ingresso >= '2025-01-01';
\end{sqlcode}

\begin{sqlcode}[caption={Exercícios 4 e 5}]
-- 4  (o curso PRIMEIRO: a chave estrangeira exige que o pai exista)
INSERT INTO curso (nome) VALUES ('Engenharia Eletrica');
INSERT INTO membro (nome, email, periodo, data_ingresso, curso_id)
VALUES ('Felipe Nunes', 'felipe@liga.edu.br', 2, '2026-03-10',
        (SELECT id FROM curso WHERE nome = 'Engenharia Eletrica'));

-- 5
SELECT * FROM membro WHERE nome = 'Carla Dias';           -- confira antes
UPDATE membro SET email = 'carla.dias@liga.edu.br' WHERE nome = 'Carla Dias';
\end{sqlcode}

\begin{sqlcode}[caption={Exercícios 6 a 8}]
-- 6  (LEFT JOIN, porque INNER JOIN esconderia quem nao tem curso)
SELECT m.nome AS membro, c.nome AS curso
FROM   membro m LEFT JOIN curso c ON m.curso_id = c.id;

-- 7
SELECT * FROM membro WHERE data_ingresso < '2024-01-01';  -- confira antes
UPDATE membro SET ativo = FALSE WHERE data_ingresso < '2024-01-01';

-- 8
INSERT INTO membro (nome, email, periodo, data_ingresso, curso_id)
VALUES ('Teste', 'teste@liga.edu.br', 1, '2026-01-01', 999);
-- ERROR 1452: a chave estrangeira exige que exista um curso com id = 999.
-- Como nao existe, o SGBD REJEITA a insercao. E a integridade referencial
-- protegendo o banco contra dado orfao.
\end{sqlcode}

\end{apendicesenv}
